% ============================================================================
% Backbone fine-tuning for appearance-invariant ISLR on ASL-Citizen.
% Self-contained methods section — \input or paste into the thesis.
% Requires amsmath.
% ============================================================================

\section{Fine-tuning the sign representation backbone}
\label{sec:backbone-ft}

\paragraph{Motivation and departure from prior approaches.}
Our earlier experiments (Sec.~\ref{tab:asl_citizen_retrieval}) trained a SignCLIP
video--text model \emph{on top of frozen} MViTv2-S ``Logos'' features: the visual backbone
was fixed and only a temporal aggregator plus a text encoder were learned with a
video--text contrastive objective. Two observations motivate a different strategy here.
First, the frozen Logos features are already highly discriminative on ASL-Citizen
(video-to-video nearest-neighbour retrieval reaches Rec@1${\approx}0.82$ with \emph{no}
training), indicating that the representational headroom lies in the \emph{backbone}, not in
a head on top of it. Second, our goal is to make the representation invariant to
\emph{appearance and signer identity}; this requires gradients to reach the visual features
themselves, which is impossible when the backbone is frozen. We therefore fine-tune the
MViTv2-S backbone directly with a sign-classification objective and an appearance-invariance
term, and evaluate the resulting \emph{frozen} features with the SignRep dictionary-retrieval
protocol (nearest-neighbour retrieval, ``Eval~A'') and a linear probe (``Eval~B'').

\paragraph{Notation.}
Let $x$ denote a $32$-frame input clip and $f_\theta$ the MViTv2-S backbone. The backbone
emits a class (\textsc{cls}) token $z = f_\theta(x) \in \mathbb{R}^{768}$, which is the global
representation used \emph{throughout} our pipeline (feature extraction, the classifier, and
all invariance objectives operate on $z$). A linear head $g_\phi$ maps $z$ to logits over the
$C{=}2731$ ASL-Citizen glosses. For a training video $i$ with gloss $y_i$ we additionally have
a set of appearance-augmented copies $\{x_i^{a}\}_{a \in A_i}$, where each augmentation
$a \in \mathcal{A} = \{\textsc{glasses},\textsc{shirt},\textsc{signer\_swap},\textsc{skin}\}$
alters appearance while preserving the sign (gloss). We write $z_i = f_\theta(x_i)$ and
$z_i^{a} = f_\theta(x_i^{a})$.

\paragraph{Base objective.}
All runs minimise a cross-entropy classification loss on the gloss labels,
\begin{equation}
\mathcal{L}_{\mathrm{CE}} = \mathrm{CE}\big(g_\phi(z),\, y\big),
\end{equation}
optionally combined with an \emph{appearance-consistency} term that pulls a video's original
representation toward those of its augmentations,
\begin{equation}
\mathcal{L}_{\mathrm{cons}} =
\frac{1}{|\mathcal{P}|}\!\!\sum_{(i,a)\in\mathcal{P}}\!\!\big(1 - \cos(z_i,\, z_i^{a})\big),
\qquad
\mathcal{L} = \mathcal{L}_{\mathrm{CE}} + \lambda\,\mathcal{L}_{\mathrm{cons}},
\label{eq:base}
\end{equation}
where $\mathcal{P}$ is the set of (original, augmentation) pairs present in the batch and
$\lambda$ weights the invariance term. This decomposes the design into three controls:
\emph{(i)}~CE on originals only (does fine-tuning the backbone help at all?), \emph{(ii)}~CE on
originals${+}$augmentations as additional data, and \emph{(iii)}~CE${+}\mathcal{L}_{\mathrm{cons}}$
(does \emph{explicitly} enforcing invariance help beyond merely seeing the augmentations?).

\subsection{Fine-tuning depth: partial, middle-ground, and full}
\label{sec:unfreeze}
MViTv2-S comprises $L{=}16$ transformer blocks in four stages. We control how much of the
backbone adapts by unfreezing only the top $k$ blocks (plus the final norm and the head),
keeping the lower $L{-}k$ blocks frozen at their pretrained Logos weights:
\begin{itemize}
  \item \textbf{Partial} ($k{=}2$): only the final stage adapts; the representation stays close
  to the pretrained features. Cheap and low-risk, but limited capacity to instil invariance.
  \item \textbf{Middle-ground} ($k{=}6$): the final stage plus part of the main stage adapt,
  with \emph{layer-wise learning-rate decay} (LLRD): block $\ell$ uses
  $\eta_\ell = \eta_0\,\gamma^{\,L-\ell}$ with decay $\gamma{=}0.75$, so deeper blocks adapt
  faster while early blocks barely move. This trades capacity against compute and the risk of
  catastrophic forgetting.
  \item \textbf{Full} ($k{=}L$): the whole backbone is fine-tuned (LLRD, low base LR). Maximum
  capacity, highest cost and forgetting risk.
\end{itemize}
This differs from the frozen-backbone SignCLIP setup, where $k{=}0$ by construction.

\subsection{Two-stage training}
\label{sec:twostage}
This recipe mirrors the schedule that gave our best frozen-feature SignCLIP result, where the
augmentation benefit only materialised when fine-tuning resumed an already-\emph{converged}
baseline rather than starting from scratch. We transpose it to the backbone:
\begin{enumerate}
  \item \textbf{Stage~1}: fine-tune from the Logos initialisation with $\mathcal{L}_{\mathrm{CE}}$
  on the original videos until convergence (validation-selected checkpoint $\theta^{(1)}$).
  \item \textbf{Stage~2}: \emph{resume} $\theta^{(1)}$ (backbone \emph{and} head), reset the
  optimiser, lower the learning rate by $100\times$, and continue with
  $\mathcal{L}_{\mathrm{CE}}+\lambda\mathcal{L}_{\mathrm{cons}}$ on originals${+}$augmentations.
\end{enumerate}
The hypothesis is that, once the classifier and features have stabilised on the in-domain gloss
distribution, the augmentations act on a settled target and the invariance term can take effect
at low LR --- whereas adding it from the Logos initialisation (our single-stage ``run~1'')
competes with the larger gradients of initial adaptation. This is the backbone analogue of the
text-anchor stabilisation argument from the SignCLIP experiments.

\subsection{Contrastive learning on the CLS token}
\label{sec:clscontrast}
The consistency term in Eq.~\eqref{eq:base} is a \emph{positive-only} pull: it aligns each
original with its augmentations but never repels different signs. Following our supervisor's
suggestion we add a proper contrastive objective on the \textsc{cls} token, in the supervised /
instance form (SupCon). Treating a video and its augmentations as a positive set
$P(u)=\{\,v \neq u : \mathrm{inst}(v)=\mathrm{inst}(u)\,\}$ (same source video) and all other
clips in the batch as negatives, with $\ell_2$-normalised tokens and temperature $\tau$,
\begin{equation}
\mathcal{L}_{\mathrm{con}}^{\textsc{cls}} =
\sum_{u}\frac{-1}{|P(u)|}\sum_{p\in P(u)}
\log\frac{\exp(z_u\!\cdot z_p/\tau)}{\sum_{v\neq u}\exp(z_u\!\cdot z_v/\tau)} .
\label{eq:supcon}
\end{equation}
Unlike $\mathcal{L}_{\mathrm{cons}}$, Eq.~\eqref{eq:supcon} includes \emph{negatives} and a
temperature, so it shapes the geometry of the $z$-space (tightening appearance-invariant
clusters while separating different videos) rather than only minimising a pairwise distance.
We emphasise that the \textsc{cls} token is already the representation our pipeline uses --- the
contribution here is the contrastive \emph{objective} applied to it, not a change of
representation.

\subsection{Output-consistency and adversarial signer-invariance (SDDA)}
\label{sec:sdda}
Following the SDDA method of \citet{fu2024sdda}, we add two further invariance terms that act on
the \emph{prediction} and through a \emph{domain adversary} rather than directly on the feature
geometry. Both operate on the same (original, augmentation) pairs $\mathcal{P}$ as
Eq.~\eqref{eq:base}, and complement the feature-space terms above.

\paragraph{Output-distribution consistency.}
Rather than pulling features together, we regularise the classifier's predictive distribution so
that an augmentation yields the same gloss posterior as its original. With
$p_i = \mathrm{softmax}(g_\phi(z_i))$ and $p_i^{a} = \mathrm{softmax}(g_\phi(z_i^{a}))$,
\begin{equation}
\mathcal{L}_{\mathrm{KL}} =
\frac{1}{|\mathcal{P}|}\!\!\sum_{(i,a)\in\mathcal{P}}\!\!\mathrm{KL}\big(p_i \,\big\|\, p_i^{a}\big).
\label{eq:sdda-kl}
\end{equation}
The SLT original sums a per-token KL over the decoded sentence; for our single-gloss classifier
this reduces to a single KL over the $C$-way gloss posterior.

\paragraph{Adversarial discrimination.}
A small domain discriminator $D$ is trained to tell originals from augmentations from the
\textsc{cls} token, while a \emph{gradient-reversal layer} (GRL, \citet{ganin2015}) trains the
backbone to \emph{fool} it, so that $z$ carries no original-vs-augmentation (appearance)
information:
\begin{equation}
\mathcal{L}_{d} = \mathrm{CE}\big(D(\mathrm{GRL}_\eta(z)),\, d\big),
\qquad d = 1\ (\text{orig}),\; 0\ (\text{aug}),
\label{eq:sdda-disc}
\end{equation}
where $\mathrm{GRL}_\eta$ is the identity in the forward pass and negates (and scales by $\eta$)
the gradient in the backward pass; $D$ is a two-layer MLP, discarded after training (the frozen
features used for evaluation are unchanged). The combined objective is
\begin{equation}
\mathcal{L} = \mathcal{L}_{\mathrm{CE}}
            + \alpha\,\mathcal{L}_{\mathrm{KL}}
            + \beta\,\mathcal{L}_{d},
\qquad \alpha = 1.0,\; \beta = 10^{-4},
\label{eq:sdda-total}
\end{equation}
matching the paper's weights. Unlike $\mathcal{L}_{\mathrm{cons}}$ (a positive-only pull on
features) and $\mathcal{L}_{\mathrm{con}}^{\textsc{cls}}$ (a contrastive reshaping of the
$z$-geometry), these act on the \emph{output} ($\mathcal{L}_{\mathrm{KL}}$) and via an
\emph{adversary} ($\mathcal{L}_d$), giving a complementary route to appearance/signer invariance;
$\mathcal{L}_{\mathrm{KL}}$ and $\mathcal{L}_d$ can also be enabled independently for ablation.

\paragraph{Evaluation.}
For every configuration we re-extract frozen $z$-features for all ASL-Citizen videos with the
fine-tuned backbone and report (A) nearest-neighbour dictionary retrieval and (B) a linear
probe, in both cases DCG, MRR and Rec@$\{1,5,10,20\}$ on the signer-disjoint test split. The
matched no-fine-tuning reference is the same backbone evaluated with identical preprocessing.

\subsection{Findings}
\label{sec:backbone-ft-findings}
Fine-tuning the backbone with cross-entropy improves over the matched no-fine-tuning reference,
confirming that the representational headroom is in the backbone. However, the
appearance-invariance interventions do \emph{not} improve in-domain ASL-Citizen recognition.
\begin{itemize}
  \item \textbf{Augmentation adds nothing.} CE on originals, CE on originals${+}$augmentations
  (aug-as-data), and the four single-augmentation consistency variants
  ($\textsc{glasses},\textsc{shirt},\textsc{signer\_swap},\textsc{skin}$) all cluster within
  ${\approx}1$ point of one another on the headline retrieval metrics (Rec@1${\approx}85$--$86$).
  This holds \emph{after} correcting a frame-rate (stride) mismatch in the
  \textsc{signer\_swap}/\textsc{skin} augmentations, ruling out a data artefact: the corrected
  variants tie with \textsc{glasses}/\textsc{shirt}.
  \item \textbf{Consistency does not help, and hurts with capacity.} At the partial depth
  ($k{=}2$) the single-stage consistency run is marginally below the CE and aug-as-data runs;
  at the middle-ground depth ($k{=}6$) it degrades clearly --- the $k{=}6$ consistency run is the
  weakest configuration overall (Eval~B Rec@1 ${\approx}82$ vs.\ ${\approx}86$ elsewhere), while
  the $k{=}6$ \emph{CE-only} run is the \emph{strongest}. The invariance term competes with the
  recognition objective, and the conflict grows as more of the backbone is allowed to move. The
  two-stage schedule (Sec.~\ref{sec:twostage}) does not recover a benefit either.
  \item \textbf{Depth, not invariance, is the lever.} The best in-domain result comes from plain
  CE with a deeper unfreeze ($k{=}6$ with LLRD), not from any appearance-augmentation strategy.
\end{itemize}
We interpret this as the backbone already being appearance-robust on ASL-Citizen: although the
test split is signer-disjoint, it shares the webcam / home-recording domain of the large
multi-dataset pretraining (Logos${+}$AUTSL${+}$WLASL), so there is little appearance \emph{shift}
for an invariance objective to exploit. A genuine test of appearance invariance therefore
requires a \emph{cross-domain} evaluation (a different capture setup and signer population); the
in-domain experiments here establish that the invariance objective is, at best, neutral, and is
harmful once it is given enough trainable capacity to compete with recognition.
